\subsection{Combat}\label{sec:qombat}

\def\NDF{\Delta}
\def\ODF{d}
\def\ODFA{d_A}
\def\ODFM{d_M}
\def\DDF{\ODF'}
\def\DDFA{\ODFA'}
\def\DDFM{\ODFM'}


Without saying any more, mechanics are already in place for combat:
it is a series of opposed actions that usually trigger rounds.
But a combat-heavy campaign could surely benefit from more detailed mechanics.
The combat module we build up in this section provides
a more detailed framework for resolving an attack.

\subsubsection{Attack Terminology}
For clarity of language we will call the acting character the \emdex{attacker}
and the opposing character the \emdex{defender}.
%There are four types of factors that determine the outcome of an attack.
We also call the aiding factors \emdex{hitting factors} and we call the difficulty factors
\emdex{preventive factors}.
These determine the \emph{success}\index{attack success} of an attack. There are two other types of factors that
contribute only to its \emph{effectiveness}\index{attack effectiveness} in the event of success:
\emdex{impact factors} and \emdex{protective factors}.
Laying out the four types of factors that contribute to attack resolution, we have:
\vspace{-1em}\begin{itemize}
\item \emph{Hitting factors} are involved in determining whether an attack hits the defender,
and they are those factors that depend on who is attacking.
\item \emph{Preventive factors} are also involved in determining whether an attack hits the defender,
but independently of who is attacking. 
\item \emph{Impact factors} contribute only to the damage a successful attack can inflict on its target
(not to its chances of success in the first place),
and they are those factors that depend on who is attacking.
\item \emph{Protective factors} also contribute only to damage inflicted,
but independently of who is attacking.
\end{itemize}
\begin{center}\begin{tabular}{l|ll}
              & offensive & defensive  \\
\cline{1-3}
success       & hitting   & prevention \\
effectiveness & impact    & protection 
\end{tabular}\end{center}
%Imagine, for example, an attacker lunging with a sword at an armored defender:
Consider the example of an attacker lunging with a sword at an armored defender.
%In the example, one hitting factor might be the sword skill of the attacker.
The attacker's sword skill is a hitting factor,
%In the example, one preventive factor might be the dodging skill of the defender.
and defender's dodging skill is a preventive factor.
Now consider what happens if the lunge strikes its target.
%In the example, one impact factor might be the strength skill of the attacker.
The attacker's strength is an impact factor,
%In the example, one protective factor might be the type of armor the defender is wearing.
and the defender's armor is a protective factor.
%It should be noted that the line between offensive and defensive factors is not always clear. Some factors will be dependent on both the attacker and the defender. For example, if an attacker sneaks up on an unsuspecting defender, the chance of a successful attack will increase; but would that be because the attacker is sneaking or because the defender is caught by surprise? Because the factors are summed together, in such a case, the element of surprise can be considered as either a hitting factor or a preventive factor; just make sure not to count its effect twice.


\subsubsection{Abilities for Combat}
There are two ways in which abilities play a role in attack resolution:
to determine whether an attack is successful
and to determine how effective it is.
When designing an ability tree, the GM
should ensure that it has an adequate set of abilities
to accommodate combat.
First, make a list of attacks that you think players could attempt in this world.
Go through all the skills and consider, for each skill, how a player might try and use the skill
in the context of an attack. Add these attacks to the list.
Now go through all the attacks in your list and ask the following of each attack:
\vspace{-1em}\begin{itemize}
\item Is it clear which ability would be relevant to determining whether the attack is successful?
\item Is it clear which ability or abilities would be relevant to determining the effectiveness of the attack?
\end{itemize}
Add or rearrange abilities until these questions can be answered satisfactorily.

\subsubsection{Resolving an Attack}
During a combat round, a player who is able to attack\index{attack} another player (e.g. is within range) may attempt to do so.
Here is the attack resolution process:
\vspace{-1em}\begin{enumerate}
\item
An attack is an opposed action, resolved according to relative performance:
$$ \rho = (R + M) - (R' + M'). $$
A negative value is a \emdex{miss}, a positive value is a \emdex{hit}, and a value of zero is a \emdex{graze}.
The modifiers $M$ and $M'$ are comprised of several bits that are explained in section \ref{sec:sitmod}.
\item
In the case of a miss, attack resolution is complete.
Otherwise determine the target of the attack, by rolling if needed (see section \ref{sec:bodyparts}).
\item
%In the case of a hit or a graze, further rolling is done to determine the severity of the outcome,
Quantify the severity of the outcome as \emdex{damage}.
Computing damage requires settling on a \emph{damage factor} $\NDF$ for the attack, explained in section \ref{sec:ndf}.
A hit results in a damage of
$$ \rho + \NDF $$
while a graze results in a damage of
$$ \frac{1}{2} \NDF, $$
rounding down.
\item
Damage then determines the type of wound\index{wound} inflicted on the target, according to section \ref{sec:wounds}.
\item
Consider what health conditions have developed, and the penalties that can influence later attacks.
\end{enumerate}

When $M$ is of a lower or greater order\index{order} than $M'$, hit damage is taken to be $0$ or $\infty$ respectively.
In this case there would be no need to consider $\NDF$.


\subsubsection{The Modifiers $M$ and $M'$}\label{sec:sitmod}


\paragraph{Situation}
Many factors having to do with the particular situation\index{situation}
surrounding an attack can be tossed into the determination of $M$ and $M'$.
Focus on
aspects of the situation that affect the attack's probability of success,
and less on the expected effectiveness of the strike in the
event of success (the latter is accounted for in impact and protection).
Here are some suggestions for situations to consider:
\vspace{-1em}\begin{itemize}
\item Weapon lengths/types. A knife-wielder is at a disadvantage defending themselves against an attacking sword-wielder.
\item Weapon condition. A properly balanced sword or a properly sighted rifle improve chances of a successful strike.
\item Cover. Someone behind cover has a defensive advantage in the face of enemy gunfire.
\item Range. A distant target is more difficult to hit with a ranged weapon.
\item Multiple attackers. Someone facing multiple attacks at the same time has more trouble defending themselves.
\item Injuries. An injured person can have some trouble defending as well as attacking.
\item Surprise. Being caught off-guard lowers one's chances of a successful defense.
\item Scale. A huge character can be much easier to hit, a small character much harder.
\item Positioning. Melee combat can be affected by things like bad footing,
lower elevation, light in eyes, kneeling, etc. %I finally found this on fudge manual p35
\end{itemize}




\paragraph{Stance}
At the beginning of each round of combat,
characters have an opportunity to choose a \emdex{stance}:
defensive ($-1$), balanced ($0$), or aggressive ($1$).
This number is added to a character's $M$ for their attack and subtracted from their
$M'$ each time they are attacked during the round.
A balanced stance is assumed if no stance is explicitly chosen.


\paragraph{Targeting}
A player may call, before their attack is resolved, that they intend to target a certain body part.
This introduces a \emdex{targeting penalty} that decreases $M$, based on the difficulty of hitting that body part.
The GM judges what the targeting penalty should be;
to give a frame of reference it could be $1$ for a torso and $4$ for an eye.
If the player does not miss, then one records in addition to the usual wounds the specific body part that received the wound.


\subsubsection{Body Parts}\label{sec:bodyparts}

%gm should create body diagram for each creature, with sizes of body parts included.
For each creature of their world, the GM should create a body diagram\index{body diagram}
that indicates relative sizes of body
parts that can serve as targets and be injured in combat.
Here is an example for humans:
%example of it for human
\begin{center} \includegraphics{svg/bp1.eps} \end{center}

%roll to determine target if no targetting was called
If no target was called by the attacking player before their attack was resolved, then a target
for the attack is chosen at random. The choice is weighted by body part size,
and it can be determined by a roll.
Mangled body parts (see section \ref{sec:wounds}) should not be included in target choice--
simply re-roll if they come up as a target.
%convert sizes to die rolls for the purpose of randomly choosing a body part, with weighting by size
Actually, the GM should convert relative body part sizes into dice roll ranges
and present the body diagram with this included.
Here is our example with d$20$ ranges put in:

%example of die roll labels for human
\begin{center} \includegraphics{svg/bp2.eps} \end{center}

\subsubsection{Impact and Protection}\label{sec:ndf}
The damage dealt by a successful attack depends on the \emdex{damage factor} $\NDF$,
the difference between \emdex{impact} and \emdex{protection}:
$$ \NDF = \text{[impact]} - \text{[protection]}. $$
%Impact is given by the impact factors, which are the factors that determine the capability of the attack to deal damage.
%The offensive damage factor $\ODF$ is defined to be the sum of the factors that determine the capability of the attack to deal damage,
%and the defensive damage factor $\DDF$ is defined to be the sum of the factors that determine the defender's capacity to absorb damage.level $\ODFA$ of a relevant ability
%plus a damage modifier $\ODFM$, and similarly for $\DDF$.
%The ability that goes into $\ODFA$ should be one that plays a role in determining how
%powerful the attack is (e.g. strength for an axe swing), and the ability that goes into $\DDFA$
%should be one that plays a role in determining a character's capacity to absorb damage (e.g. constitution). 
%If no relevant ability can be identified for $\ODFA$ or $\DDFA$, set it to $0$. 
%The damage modifer consists of other factors that affect the ability of the attack to do damage.
%For example $\ODFM$ could involve the weapon being used, and $\DDFM$ could involve the defender's armor.
%Putting it together,
%$$ \NDF = (\ODFA + \ODFM) - (\DDFA + \DDFM). $$
Impact is a sum of effects coming from the impact factors involved in the attack,
and protection is a sum of effects coming from protection factors.
Let us go through some of the factors that should be included.

\paragraph{Abilities}
The level of a relevant ability can be included directly in impact.
Or, if it only partly applies, a lesser portion of the ability level may be used.
For example a stronger character deals a deadlier blow with a melee weapon, so a strength ability could be used.
Don't stack the effects of multiple abilities; instead, use the nearest common ancestor attribute when multiple abilities apply.
A relevant ability may also be included, separately, in protection.
For example a tougher character can take a harder hit without receiving as much of an injury.

\paragraph{Size}
The same blow imparts greater destruction to smaller targets. 
This is why the \emdex{size} of a target is important to protection.
When creating their world, the GM should choose a \emdex{baseline target}
size with respect to which all target sizes are compared.
The choice of baseline target calibrates the entire damage system to respond appropriately to target sizes.
A good baseline target is one that, with all other things being equal, would be injured a bit from a good punch.
More precisely, it is a target that recieves a wound of degree $1$ (see section \ref{sec:wounds}) about $38\%$
of the time, and a wound of degree $0$ the rest of the time, from a successful unarmed attack by an equally matched standard being.
%use the protection table to determine suitable contribution to protection due to size of body part
Use the following table to choose a suitable contribution to protection due to the size of a body part:
\begin{center}
\input{size_protection_table.tex}
\end{center}
Sizes are measured in in units of the baseline target size,
and the value in the right hand column is included in the calculation of total protection.
%we recommend just putting this number on the body charts directly
Actually, the GM can include these protection modifiers on the body diagram\index{body diagram}
directly for players' convenience.
%example for human, with torso as baseline target
Here is our human example from above, with protection modifiers included,
using a human torso as the baseline target:

\begin{center} \includegraphics{svg/bp3.eps} \end{center}

\paragraph{Other Things}
We list below suggestions for some other factors that can be included.
An understanding of section \ref{sec:wounds} is required for the GM to properly
judge the contribution of these factors to impact and protection.
\vspace{-1em}\begin{itemize}
\item Weapon quality. A sharp sword has a greater impact than a dull or worn sword.
\item Weapon type. A knife has greater impact than a fist.
As a baseline for this modifier, a contribution of $0$ to impact
can be attributed to an unarmed attack by the standard being.
\item Armor, natural or otherwise. A creature with a tough exterior or armored clothing has greater protection.
\end{itemize}


\subsubsection{Wounds}\label{sec:wounds}

%intro
%hp of body parts
Wounds on body parts are tracked in terms of \emdex{hit points} (hp).
Each body part of each character has a number of hit points ranging from
$0$ to $4$.
%wound of 'degree' x
When damage dealt to a target body part results in a decrease of its hp by an amount $n$, 
a \emph{wound of degree}\index{wound degree}
$n$ is said to have been inflicted upon the target.
%damage dealt is according to following table. record new hp. on body diagram, why not.
Use the following table to determine the degree of wound inflicted on a defender's body part
due to a given damage:
\begin{center}
\input{damage_table.tex}
\end{center}
Players should keep track of the hp values of body parts of the characters under their control;
the body diagram is a useful place to record this information.

%pay attention to health peups that could have been triggered
Some health conditions may have developed following a wound;
remember to consider how resulting penalties can influence later attacks.

%death
The defender dies if a \emdex{vital body part} is at an hp of $0$.
The GM should decide, for each creature of their world, which body parts are vital.

%narrative interpretation of wounds
The hp values for a body part have the following descriptions that can feed into the narrative:
\begin{center}
{\setlength{\extrarowheight}{1pt}
\begin{tabular}{l|l}
hp & adjective \\
\cline{1-2}
4  & fine \\
3  & scratched \\
2  & hurt \\
1  & incapacitated \\
0  & mangled
\end{tabular}}
\end{center}

\paragraph{Healing} Wounds are treated like any other health condition, so healing\index{healing} is governed by a flowchart.
After deciding on the creatures that inhabit their world, the GM should devise for each one a healing flowchart.
Here is a somewhat generic healing flowchart geared towards organic creatures:
\begin{center} \includegraphics{svg/healing.eps} \end{center}
